% !PRINTED 101
\weberpage{101}
%<<CONT
different, and cannot be carried over into one another by \gesperrt{Tschirnhausen transformation}, because they belong to primary divisors of different full cyclotomic fields, or, if they are contained in the same full cyclotomic field, have different groups.

One can, to be sure, also form cyclotomic periods from other roots of unity, which lead to cubic equations. These periods, however, can be expressed by lower roots of unity, and are not primary. Thus one obtains, e.g., when $r$ is a 35th root of unity, a period of 8 terms whose exponents are formed from the powers of $-8$:

\[
\eta = r + r^{-1} + r^6 + r^{-6} + r^8 + r^{-8} + r^{13} + r^{-13},
\]
which therefore also satisfies a cubic equation.

But we have
\[
1 + r^7 + r^{-7} + r^{14} + r^{-14} = 0,
\]
and if one multiplies this equation by $(r^{15} + r^{-15})$:
\[
r^{15} + r^{-15} + r + r^{-1} + r^6 + r^{-6} + r^8 + r^{-8} + r^{13} + r^{-13} = 0,
\]
hence
\[
\eta = -(r^{15} + r^{-15}),
\]
which already occurs among the periods of the 7th roots of unity. This phenomenon is explained by the fact that $m = 35$ does not have the property required in \S.\,21, that $q - 1$ is divisible by 3 for all primes $q$ dividing $m$.

\begin{center}
\S.\,24.

\gesperrt{Biquadratic Cyclotomic Fields.}
\end{center}

With the biquadratic cyclotomic equations one has two kinds to distinguish. We can assume two invariants $i_1 = i_2 = 2$ or only one invariant $i_1 = 4$. We shall put the first case first.

For $m$ we have in this case, according to the summary in \S.\,21, the following conditions:

$m$ may be odd, or divisible by 4 or by 8, but by no higher power of 2. No odd prime can divide $m$ more than once, and there must be at least two index moduli present; hence, besides in
%CONT>>
